% Generated draft: requires mathematical and source review.
\documentclass[11pt]{article}
\usepackage[margin=25mm]{geometry}
\usepackage{fontspec}
\setmainfont{Latin Modern Roman}
\usepackage{amsmath,amssymb,mathtools,unicode-math}
\setmathfont{Latin Modern Math}
\usepackage{enumitem,needspace,xurl,hyperref}
\setlength{\parindent}{0pt}
\setlength{\parskip}{5pt}
\setlength{\emergencystretch}{3em}
\newcommand{\N}{\mathbb{N}}
\newcommand{\Z}{\mathbb{Z}}
\newcommand{\Q}{\mathbb{Q}}
\newcommand{\R}{\mathbb{R}}
\newcommand{\C}{\mathbb{C}}
\newcommand{\F}{\mathbb{F}}
\newcommand{\A}{\mathbb{A}}
\newcommand{\PP}{\mathbb{P}}
\newcommand{\E}{\mathbb{E}}
\newcommand{\eps}{\varepsilon}
\newcommand{\dd}{\,\mathrm{d}}
\begin{document}
\section*{UnsolvedMath ID: 30006012}
\textbf{Source problem code:} catalogue entry
\subsection{Definitions and mathematical statement}
Let \(\A_f\) denote the finite adeles of \(\Q\), and let \(K\subset\operatorname{GSp}_4(\A_f)\) be an open compact subgroup. The supplied source describes \(S_K(\operatorname{GSp}_4)\) as a three-dimensional Shimura variety parametrizing principally polarized abelian surfaces with level structure determined by \(K\). No additional conditions on \(K\) are supplied. For a prime \(\ell\), write
\[
V_K^3=H^3_{\mathrm{\acute et}}\!\left(S_K(\operatorname{GSp}_4)_{\overline{\Q}},\Q_\ell\right)
\]
for the geometric degree-three cohomology suggested by the record; the original excerpt does not unambiguously preserve the geometric base notation.

The motivating context is a classical cuspidal eigenform \(g\) of weight two for \(\operatorname{GL}_2\). The source attributes to Deligne a Galois representation \(\rho_g\), described as a quotient of the degree-one étale cohomology of a modular curve \(X_1(N)\) for sufficiently large \(N\). It states that \(\rho_g\) also occurs in the cohomology of other Shimura varieties, including the indicated \(\operatorname{GSp}_4\) setting. The coefficient field, any required extension of \(\Q_\ell\), and the precise occurrence are not specified. The excerpt begins a decomposition involving \(\Pi_f^K\otimes\rho_{\Pi_f}\), but does not supply a complete formula, its indexing set, or the definitions and hypotheses for these factors.

For comparison, in the usual cycle-class formulation, if \(X\) is a smooth variety over a number field \(k\) and \(r\geq 0\) is an integer, codimension-\(r\) algebraic cycles have classes under
\[
\operatorname{cl}_\ell:\operatorname{CH}^r(X_{\overline{k}})\otimes_{\Z}\Q_\ell
\longrightarrow H^{2r}_{\mathrm{\acute et}}(X_{\overline{k}},\Q_\ell(r)).
\]
Here \(\operatorname{CH}^r\) denotes cycles modulo rational equivalence, and \((r)\) denotes the Tate twist. A Tate class in this target is a class fixed by \(\operatorname{Gal}(\overline{k}/L)\) for some finite extension \(L/k\).

\textbf{Source question, with unresolved target.} Determine when the classes referred to in the source, extracted from the relevant degree-three cohomology of \(S_K(\operatorname{GSp}_4)\), arise from algebraic cycles in the intended Tate-conjecture construction. The record does not identify these extracted classes or the cycle-class map intended to produce them. Since ordinary codimension-\(r\) cycle classes lie in degree \(2r\), they do not directly lie in \(V_K^3\). A product with another variety, a correspondence, or another cohomological construction might be intended, but none is specified. Consequently, a precise algebraicity or surjectivity assertion cannot be recovered without further source information; no such construction is assumed here.
\subsection{Short English statement}
The source asks when classes associated with the degree-three cohomology of the \(\operatorname{GSp}_4\) Shimura threefold come from algebraic cycles. Its excerpt ends before identifying the relevant classes and construction. Ordinary cycle classes have even cohomological degree, so the missing construction is essential to a precise statement. The supplied open status is dated curation, not an independently verified conclusion.
\subsection{Sources}
\begin{itemize}
\item Local UnsolvedMath registry, supplied complete record: permanent ID \(30006012\), source problem code \texttt{catalogue entry}, title \emph{Algebraicity of Tate Classes on GSp4 Shimura Varieties}, in the Oberwolfach Reports Open Problems collection. No proposer is named. The record supplies proposed year \(2024\), status \texttt{open}, and editorial statement status \texttt{corrected\_verified}; its verification note says only formula delimiters, notation, grammar, and typography were repaired while preserving mathematical content and scope. These labels do not establish mathematical correctness.
\item Supplied report citation: \emph{Arithmetic Geometry (2025)}, Oberwolfach Reports, DOI \texttt{10.4171/owr/2024/33}, pp.~1855--1912. \url{https://doi.org/10.4171/owr/2024/33}. The background describes these reports as records of workshops and their open-problem discussions. The citation is retained as supplied; the report title's year and the DOI/proposed-year information require reconciliation.
\item Supplied literature source labeled ``Sweeting'': \url{https://arxiv.org/abs/2211.10838}. The record supplies no further bibliographic details or exact result identifying the intended classes. The original excerpt attributes the construction of \(\rho_g\) to Deligne, without a separate citation.
\item Literature checked at: \texttt{2026-08-22}. This is supplied dated curation, not independently verified. The supplied assessment reports substantial results for endoscopic and generic cases, while stating that no general algebraicity theorem for all extracted classes is known. Its listed sources are the OWR report and Sweeting link above. The background explicitly calls this dated literature triage derived from the supplied OWR curation and requests independent verification of time-sensitive status and cited claims.
\end{itemize}
\end{document}
